
\par\noindent\textit{引理~\ref{lem:7-2-291} 的证明（续）.}\quad
下面证明 (7.2.292.4)。
$B$ 的满射性由易证的 (7.2.137.2) 得出；
$B\circ A=0$ 则由 (7.2.292.2) 得出。
设
\[
\mathfrak C_{r=r_0}
=\sum_{a\in I_{r=r_0}}c_a^{r=r_0}\mathfrak P_a^{r=r_0},
\qquad
\mathfrak C_{s=s_0}
=\sum_{a\in I_{s=s_0}}c_a^{s=s_0}\mathfrak P_a^{s=s_0},
\]
并设
$B(\mathfrak C_{s=0},\mathfrak C_{s=1},
\mathfrak C_{r=0},\mathfrak C_{r=1})=0$。
以下四元组先列 $s=0,1$ 端面，再列 $r=0,1$ 端面。
我们要证明
\[
\mathbf C:=(\mathfrak C_{s=0},\mathfrak C_{s=1},
\mathfrak C_{r=0},\mathfrak C_{r=1})\in\operatorname{Im}A.
\]

令 $I_{r=0,(0,0)}$ 是所有 $a\in I_{r=0}$ 中
$\mathfrak P_a^{r=0}$ 与
$\{(0,0)\}\times L'_1$ 相交者组成的集合；
同理定义 $I_{s=0,(0,0)}$。
置
\[
\begin{aligned}
\mathfrak C_{r=0,(0,0)}
&=\sum_{a\in I_{r=0,(0,0)}}
c_a^{r=0}\mathfrak P_a^{r=0},\\
\mathfrak C_{s=0,(0,0)}
&=\sum_{a\in I_{s=0,(0,0)}}
c_a^{s=0}\mathfrak P_a^{s=0}.
\end{aligned}
\]
由假设，
\[
\operatorname{Eval}_{s=0}(\mathfrak C_{r=0,(0,0)})
=\operatorname{Eval}_{r=0}(\mathfrak C_{s=0,(0,0)}).
\]
记此公共链为 $\mathfrak Y$。

对 $a\in I_{r=0,(0,0)}$，定义
\[
\widetilde f_a\colon
|\mathfrak P_a^{r=0}|\times[0,1]
\longrightarrow[0,1]^2\times L'_1,
\qquad
\widetilde f_a(x,r)=(r/4,f_a(x)),
\]
其中在右端把
$f_a(x)\in\{0\}\times[0,1]_s\times L'_1$
识别为 $[0,1]_s\times L'_1$ 中的点。
对 $(|\mathfrak P_a^{r=0}|\times[0,1],\widetilde f_a)$
取适当的棱柱剖分，得到
$\phi(\mathfrak P_a^{r=0})
\in S_+([0,1]^2\times L'_1)$，满足
\[
\begin{aligned}
\operatorname{Eval}_{r=0}
\bigl(\phi(\mathfrak P_a^{r=0})\bigr)
&=\mathfrak P_a^{r=0},&
\operatorname{Eval}_{s=0}
\bigl(\phi(\mathfrak P_a^{r=0})\bigr)
&=\phi\bigl(\operatorname{Eval}_{s=0}
(\mathfrak P_a^{r=0})\bigr),\\
\operatorname{Eval}_{s=1}
\bigl(\phi(\mathfrak P_a^{r=0})\bigr)
&=0,&
\operatorname{Eval}_{r=1}
\bigl(\phi(\mathfrak P_a^{r=0})\bigr)
&=0.
\end{aligned}
\]
取下列双参数奇异链：
\[
\mathfrak E_1=
\sum_{a\in I_{r=0,(0,0)}}
c_a^{r=0}\phi(\mathfrak P_a^{r=0})
\in S_+([0,1]^2\times L'_1).
\]
其端面像满足
\[
A(\mathfrak E_1)
=(\phi(\mathfrak Y),0,\mathfrak C_{r=0,(0,0)},0).
\]
从 $\mathbf C$ 中扣除 $A(\mathfrak E_1)$，
便可化为 $\mathfrak C_{r=0,(0,0)}=0$ 的情形。
对于修正后的四元组，
(7.2.292.2) 又使其角点公共链为零。
交换 $s$ 与 $r$ 重复以上论证，得到 $\mathfrak E_2$；
再扣除 $A(\mathfrak E_2)$，
可使 $\mathfrak C_{r=0,(0,0)}
=\mathfrak C_{s=0,(0,0)}=0$。
也就是说，经两次修正后，四元组不再与
$\{(0,0)\}\times L'_1$ 相交。
记这一角点的累计修正链为\translatornote{原书先后从同一四元组扣除
$A(\mathfrak E_1)$ 和 $A(\mathfrak E_2)$，
下句却把总修正链写成 $\mathfrak E^{(0,0)}=\mathfrak E_2$；
依加法性应为 $\mathfrak E_1+\mathfrak E_2$。}
\[
\mathfrak E^{(0,0)}
:=\mathfrak E_1+\mathfrak E_2
\in S_+([0,1]^2\times L'_1).
\]

在其余三个角点重复此法，得到
$\mathfrak E^{(1,0)}$、$\mathfrak E^{(0,1)}$、
$\mathfrak E^{(1,1)}$，使
\[
\begin{aligned}
&\mathbf C-A\bigl(
\mathfrak E^{(0,0)}+\mathfrak E^{(1,0)}
+\mathfrak E^{(0,1)}+\mathfrak E^{(1,1)}\bigr)\\
&\hspace{3em}=:
(\mathfrak C'_{s=0},\mathfrak C'_{s=1},
\mathfrak C'_{r=0},\mathfrak C'_{r=1})
\end{aligned}
\]
的四个分量都不与角点相交。

将不与角点相交的 $r=0$ 分量写成
\[
\mathfrak C'_{r=0}
=\sum_{a\in A}c_a\mathfrak P_a^{r=0}.
\]
由于每个 $\mathfrak P_a^{r=0}$ 都不与角点相交，
有
\[
A\Bigl(\sum_{a\in A}c_a\phi(\mathfrak P_a^{r=0})\Bigr)
=(0,0,\mathfrak C'_{r=0},0).
\]
对 $r=1$ 及 $s=0,1$ 重复相同构造，
可知上述四元组属于 $\operatorname{Im}A$。
引理~\ref{lem:7-2-291} 得证。$\square$
